\documentclass[12pt]{article}
\usepackage[utf8]{inputenc}
\usepackage[T1]{fontenc}
\usepackage{hyperref}
\title{Calamian Tagbanwa Language Resources for K-12 Curriculum Development}
\author{Compiled by Claude Sonnet 5 (claude-sonnet-5) for Andr\'as Kornai}
\date{}
\begin{document}
\maketitle
\section{Introduction}
Calamian Tagbanwa (ISO 639-3 tbk; Glottocode cala1258; Austronesian) is spoken in the Calamian Islands north of Palawan, Philippines (Busuanga, Coron, Culion, Linapacan). Estimates are about 10,000 speakers (2007), though a Webonary page gives about 25,000 people. It is written in Latin script and, to some extent, the Tagbanwa script. Search results gave little on dictionaries or schooling.
\section{Primary Educational Resources}
No primers or school programmes were found. Webonary states that literature has been produced and that mother-tongue literacy is significant, but no specific titles were seen.
\section{Linguistic and Reference Materials}
\begin{itemize}
\item Webonary Tagbanwa dictionary project and introduction. \url{https://www.webonary.work/tagbanwa/overview/introduction/?lang=en} \quad \url{https://www.webonary.work/tagbanwa/?p=9444}
\item Omniglot page on Calamian Tagbanwa and its scripts. \url{https://www.omniglot.com/writing/calamiantagbanwa.htm}
\item Wikipedia article. \url{https://en.wikipedia.org/wiki/Calamian_Tagbanwa_language}
\end{itemize}
\section{Religious and Community Resources}
Joshua Project lists a Baras Kalamian Tagbanwa people-group page; no scripture or community organisation was verified. \url{https://joshuaproject.net/people_groups/12445}
\section{Digital and Audio Resources}
A verb list on Verbix. \url{https://docs.verbix.com/Languages/TagbanwaCalamian}
\section{Curriculum Development Context}
Inter-marriage with non-Tagbanwa speakers is reducing child use. The Webonary dictionary and the Tagbanwa script are the main starting points; community and DepEd mother-tongue programme contacts were not found and should be sought directly.
\end{document}